% ID: MAT-TRI-EQG-018
% Titolo: Intersezioni di una funzione goniometrica con gli assi (1)
% Disciplina: Matematica
% Area: trigonometria
% Argomento: equazioni_goniometriche
% Classe: Quarta secondaria di secondo grado
% Difficolta: 3
% Tipo: grafico
% Tag: matematica, trigonometria, equazioni_goniometriche
% Fonte: verifica_goniometria_base_anonimizzata
% Provenienza: materiale fornito dal docente; testo anonimizzato; soluzione redatta per Exergo

\begin{esercizio}
Trova le coordinate dei punti di intersezione con gli assi
cartesiani del grafico della funzione: 
\[
f(x)=2\sin^2 x+4\cos^2 x
-2\sqrt{3}\sin x\cos x-2.
\]
Considera $x\in[0;2\pi)$.
Salvo diversa indicazione, lavora in $\mathbb R$; angoli in radianti.
\end{esercizio}

\begin{soluzione}
Usando $\sin^2x+\cos^2x=1$ si ottiene
$f(x)=2\cos x(\cos x-\sqrt3\sin x)$.
Gli zeri nell'intervallo sono $\pi/6,\pi/2,7\pi/6,3\pi/2$.
Le intersezioni sono $(a,0)$ per queste ascisse e $(0,2)$.
\end{soluzione}
